\chapter{Adatbázis tervezés és felépítés}

\section{Globális megközelítés}
A perzisztencia réteg célja a felhasználók, a felhasználóhoz kötött szó- és színpaletták, valamint a táblaállapot-mentések tárolása \emph{referenciális integritással}. A tipikus éles környezet \textbf{InnoDB} motort és \textbf{utf8mb4} karakterkészletet használ; fejlesztéshez gyakori az SQLite. A részletes \texttt{CREATE TABLE} szövegek a \texttt{api/docs/database-schema.md} fájlban találhatók — itt a logikai szerkezet és a kulcsok összefoglalója szerepel.

\section{Üzleti táblák áttekintése}
\begin{longtable}{@{}p{3.2cm}p{10.5cm}@{}}
\toprule
\textbf{Tábla} & \textbf{Szerep és fő mezők} \\
\midrule
\endfirsthead
\multicolumn{2}{c}{\tablename\ \thetable{} — folytatás} \\
\toprule
\textbf{Tábla} & \textbf{Szerep és fő mezők} \\
\midrule
\endhead
\midrule
\multicolumn{2}{r}{Folytatódik a következő oldalon\ldots} \\
\endfoot
\bottomrule
\endlastfoot

\texttt{users} &
Elsődleges kulcs \texttt{id}. Egyedi \texttt{email} és \texttt{username}. Mezők: \texttt{name}, \texttt{password} (hash), \texttt{role} (pl. \texttt{vendeg}, \texttt{diak}, \texttt{tanar}, \texttt{admin}), \texttt{active}, \texttt{suspended\_at}, \texttt{email\_verified\_at}, időbélyegek. A bejelentkezés tiltása felfüggesztett vagy inaktív állapotban történik. \\

\texttt{lists} &
Felhasználói szólisták konténere: \texttt{user\_id} $\rightarrow$ \texttt{users.id}, \texttt{name}. Egy felhasználóhoz több lista (1:N). \\

\texttt{words} &
Listaelemek: \texttt{list\_id} $\rightarrow$ \texttt{lists.id}, \texttt{word} (max 255), \texttt{position} ($\geq 0$). \textbf{Egyediség:} \texttt{UNIQUE(list\_id, word)} és \texttt{UNIQUE(list\_id, position)} — ez indokolja az admin kliens kétlépcsős pozíció-frissítését cserekor. \\

\texttt{color\_lists} &
Színpaletta-listák: \texttt{user\_id}, \texttt{name}. \\

\texttt{colors} &
Egy paletta színei: \texttt{list\_id} $\rightarrow$ \texttt{color\_lists.id}, \texttt{color} (string, tipikusan hex), \texttt{position}. \textbf{Egyediség:} \texttt{UNIQUE(list\_id, position)}. \\

\texttt{board\_save\_groups} &
Mentési mappák: \texttt{user\_id} $\rightarrow$ \texttt{users.id} (\textbf{ON DELETE CASCADE}), \texttt{name}, opcionális \texttt{position} a megjelenítési sorrendhez. \\

\texttt{board\_saves} &
Konkrét mentés: \texttt{board\_save\_group\_id} $\rightarrow$ \texttt{board\_save\_groups.id} (\textbf{CASCADE}), \texttt{user\_id} $\rightarrow$ \texttt{users.id} (\textbf{CASCADE}), \texttt{name}, \texttt{payload} (JSON, NOT NULL). \textbf{Egyediség:} \texttt{UNIQUE(board\_save\_group\_id, name)}. \\
\end{longtable}

\section{ER és UML}
A \texttt{api/docs/database-uml.md} Mermaid és PlantUML diagramokat tartalmaz; a lényegi üzenet: a \texttt{User} központi entitás, amelyhez a listák, színes listák és mentési csoportok csatlakoznak; a \texttt{BoardSave} kettős kapcsolatot mutat (user + csoport) a gyors jogosultság-ellenőrzés és az integritás miatt.

\section{Laravel infrastruktúra táblák}
Az alkalmazás működéséhez további táblák szükségesek (nem üzleti domain):
\texttt{migrations}, \texttt{sessions}, \texttt{cache}, \texttt{cache\_locks}, \texttt{jobs}, \texttt{job\_batches}, \texttt{failed\_jobs}, \texttt{password\_reset\_tokens}, \texttt{personal\_access\_tokens}.
Utóbbi tárolja a Sanctum API tokeneket; enélkül a Bearer hitelesítés nem működik.

\section{Nem funkcionális szempontok (adat)}
\begin{itemize}
  \item A \texttt{board\_saves.payload} nagy lehet — lista lekérdezéskor a jelen backend a teljes payloadot is visszaadja minden mentéshez (lásd követelmény-specifikáció az eltérésről az \texttt{api-board-saves.md} javaslatához képest).
  \item JSON típus: MySQL/MariaDB natív JSON; SQLite fejlesztői módban Laravel kompatibilitást biztosít.
\end{itemize}
